\section{Method}
\label{sec:method}

The system has three parts: a per-review summarizer, a map--reduce procedure
that extends it to many reviews, and an aspect miner that answers the question
the summarizer structurally cannot.

\subsection{Per-review summarization}

We fine-tune \texttt{t5-small} \cite{raffel2020t5} on (review, headline) pairs.
T5 is a multi-task text-to-text model in which the task is selected by a
textual prefix, so each input is prepended with \texttt{"summarize: "}.
Decoding uses beam search with four beams and blocks repeated bigrams; without
the latter the model degenerates into repetition (\emph{``good good good''}) on
targets this short.

\subsection{Map--reduce aggregation}

A product page carries tens or hundreds of reviews, and the encoder accepts at
most 256 tokens. Concatenating the reviews and passing them in one call does
not fail --- it silently truncates, so most reviews never influence the output
at all. We therefore aggregate in two stages, following the map--reduce shape:
each review is summarised independently (\emph{map}), the resulting summaries
are joined into a single pseudo-document, and that document is passed back
through the same model (\emph{reduce}) with a doubled generation budget.

This raises the ceiling but does not remove it, and we state the consequence
plainly: the reduce step is itself subject to the 256-token encoder limit. With
per-review summaries averaging four to six tokens, a single reduce pass
saturates somewhere around forty to fifty reviews, beyond which the later
summaries are truncated exactly as the raw reviews would have been. A
hierarchical reduce --- reducing in batches, then reducing the results ---
would remove the ceiling entirely. We did not implement it, and the
demonstration in Section~\ref{sec:analysis} uses sixty reviews, which is inside
the regime where mild truncation is possible. This is a limitation of our
implementation, not of the approach.

\subsection{Aspect mining}

Neither a per-review summary nor a reduced summary answers the question the
task actually poses --- which aspects do customers as a group praise, and which
do they criticise. That answer is a property of the collection, not of any
review in it. We compute it separately.

Reviews are split into sentences and each sentence is scored with VADER
\cite{hutto2014vader}, a rule-based sentiment analyser. Sentences above $+0.2$
form the positive pool and those below $-0.2$ the negative pool; the stricter
sentence threshold means only clearly polar sentences contribute aspects. A
separate, looser threshold of $\pm0.05$ on the whole review is used for the
positive/negative/neutral counts reported to the user.

Candidate aspects are unigrams and bigrams, with English stop words removed
along with a hand-built list of \emph{opinion} words. This second list matters:
terms such as \emph{good}, \emph{love} and \emph{disappointed} carry the
verdict rather than the topic, and terms such as \emph{product},
\emph{amazon} and \emph{bought} are true of every review. Neither identifies
what is being discussed.

Ranking candidates by raw frequency does not work. For a cat toy, \emph{cat}
and \emph{toy} head both lists, and praise and criticism look identical. We
instead select \emph{contrastively}: a term is listed on a side only if the
share of its usage falling on that side, measured as a rate rather than a raw
count because the two pools differ in size, is at least $0.6$ --- that is,
mentioned at least $1.5\times$ more often there. Terms with fewer than two
mentions are dropped as noise. Finally, a unigram is discarded when it is
dominated by a bigram containing it, so \emph{battery life} is preferred to
\emph{battery} while a standalone \emph{flavor} survives.

We emphasise what this is not. It is a frequency-and-polarity heuristic, not a
trained aspect-based sentiment model. It identifies \emph{what} is discussed on
each side; it cannot resolve a negated aspect inside an otherwise positive
sentence.

\subsection{The reported verdict}

The paragraph the application shows is assembled from three sources rather
than generated end to end, and we describe it as such. A mood phrase is chosen
from the share of \emph{polar} reviews that are positive --- neutral reviews
are excluded from this ratio deliberately, since counting them drags a
genuinely mixed product toward ``unhappy'' --- followed by the top five
praised and criticised aspects, followed by the reduce step's generated
sentence quoted verbatim. The neural component contributes the summaries and
the quoted verdict; the structure around it is a template.
